\chapter{2025年天津卷}

\section{单选题}

\begin{problem}
已知集合 \(\mathcal U=\{1,2,3,4,5\}\)，\(A=\{1,3\}\)，\(B=\{2,3,5\}\)，则 \(\mathcal U\smallsetminus(A\cup B)=\)（\quad）

\choices
    {\(\{1,2,3,4\}\)}
    {\(\{2,3,4\}\)}
    {\(\{2,4\}\)}
    {\(\{4\}\)}
\end{problem}

\begin{answer}
D
\end{answer}

\begin{solution}
按本页现稿与原PDF题面，计算结果均为 \(\{4\}\)，对应现稿 D与原PDF D，故结论一致. 因为 \(A=\{1,3\}\)，\(B=\{2,3,5\}\)，所以 \(A\cup B=\{1,2,3,5\}\)，故 \(\mathcal U\smallsetminus\left( A\cup B \right)=\{4\}\). 对照现稿选项，只有 D为 \(\{4\}\)，故排除 A，B，C；对照原PDF选项，同样只有 D为 \(\{4\}\)，故排除 A，B，C.
\end{solution}

\begin{problem}
设 \(x\in\R\)，则“\(x=0\)”是“\(\sin 2x=0\)”的（\quad）

\choices
    {充分不必要条件}
    {必要不充分条件}
    {充分必要条件}
    {既不充分也不必要条件}
\end{problem}

\begin{answer}
A
\end{answer}

\begin{solution}
按本页现稿与原PDF题面表述等价，结论均为 A. 因为 \(x=0\) 时 \(\sin 2x=\sin 0=0\)，所以“\(x=0\)”是“\(\sin 2x=0\)”的充分条件；因为存在 \(x=\pi\) 使 \(\sin 2x=0\) 但 \(x\ne 0\)，所以不是必要条件. 故是充分不必要条件，选 A，故排除 B，C，D，其中 B颠倒了方向，C要求双向成立，D否认了已成立的充分性.
\end{solution}

\begin{problem}
已知函数 \(y=f(x)\) 的图象如下，则 \(f(x)\) 的解析式可能为（\quad）

\bitmapfigure[max width=0.42\linewidth]{img/2025/tianjin/q3_fig1.png}

\choices
    {\(f(x)=\frac{x}{1-|x|}\)}
    {\(f(x)=\frac{x}{|x|-1}\)}
    {\(f(x)=\frac{|x|}{1-x^2}\)}
    {\(f(x)=\frac{|x|}{x^2-1}\)}
\end{problem}

\begin{answer}
D
\end{answer}

\begin{solution}
按本页现稿与原PDF题面条件等价，结论均为 D. 因为图象关于 \(y\) 轴对称，所以 \(f\left( x \right)\) 为偶函数. 因为选项 A中 \(f\left( -x \right)=-f\left( x \right)\)，选项 B中 \(f\left( -x \right)=-f\left( x \right)\)，所以 A，B均为奇函数，故排除 A，B. 因为当 \(0<x<1\) 时 \(\frac{\left| x \right|}{1-x^{2}}>0\)，\(\frac{\left| x \right|}{x^{2}-1}<0\)，而图象在 \(\left( 0,1 \right)\) 内为负，所以排除 C. 因为当 \(x>1\) 时选项 A，C均为负，与右支为正不符，而选项 D在 \(\left( 0,1 \right)\) 内为负，在 \(\left( 1,+\infty \right)\) 内为正且趋于 \(0\)，与图象相符，故选 D.
\end{solution}

\begin{problem}
若 \(m\) 为直线，\(\alpha\)，\(\beta\) 为两个平面，则下列结论中正确的是（\quad）

\choices
    {若 \(m\parallel \alpha\)，\(n\subset \alpha\)，则 \(m\parallel n\)}
    {若 \(m\perp\alpha\)，\(m\perp\beta\)，则 \(\alpha\perp\beta\)}
    {若 \(m\parallel\alpha\)，\(m\perp\beta\)，则 \(\alpha\perp\beta\)}
    {若 \(m\subset\alpha\)，\(\alpha\perp\beta\)，则 \(m\perp\beta\)}
\end{problem}

\begin{answer}
C
\end{answer}

\begin{solution}
按本页现稿 A项与原PDF A项不同，但结论均为 C. 现稿 A为“若 \(m\parallel \alpha\)，\(n\subset \alpha\)，则 \(m\parallel n\)”，原PDF A为“若 \(m\parallel \alpha\)，\(m\parallel \beta\)，则 \(\alpha\parallel \beta\)”. 因为现稿 A中 \(m\) 与 \(n\) 可为异面直线，所以现稿 A错误；因为存在两个相交平面同时平行于同一条直线，所以原PDF A错误，故两版 A均排除. 因为若 \(m\perp \alpha\) 且 \(m\perp \beta\)，则 \(\alpha\) 与 \(\beta\) 的法向量平行，所以 \(\alpha\parallel \beta\)，而不是 \(\alpha\perp \beta\)，故 B错误. 因为 \(m\parallel \alpha\) 时在 \(\alpha\) 内存在直线 \(n\) 使 \(n\parallel m\)，又 \(m\perp \beta\)，所以 \(n\perp \beta\)，而 \(\alpha\) 内存在一条直线垂直于 \(\beta\)，所以 \(\alpha\perp \beta\)，故 C正确. 因为 \(m\subset \alpha\) 且 \(\alpha\perp \beta\) 时 \(m\) 可与 \(\beta\) 平行，相交或在 \(\beta\) 内，不一定垂直，所以 D错误. 故选 C.
\end{solution}

\begin{problem}
下列说法中错误的是（\quad）

\choices
    {若 \(X\sim N(\mu,\sigma^2)\)，则 \(P(X\le\mu-\sigma)=P(X\ge\mu+\sigma)\)}
    {若 \(X\sim N(1,2^2)\)，\(Y\sim N(2,2^2)\)，则 \(P(X<1)<P(Y<2)\)}
    {\(|r|\) 越接近 1，相关性越强}
    {\(|r|\) 越接近 0，相关性越弱}
\end{problem}

\begin{answer}
B
\end{answer}

\begin{solution}
按本页现稿与原PDF题面条件等价，结论均为 B. 因为 \(X\sim N\left( \mu,\sigma^{2} \right)\) 关于 \(X=\mu\) 对称，所以 \(P\left( X\le \mu-\sigma \right)=P\left( X\ge \mu+\sigma \right)\)，故 A正确. 因为 \(X\sim N\left( 1,2^{2} \right)\) 时 \(\frac{X-1}{2}\sim N\left( 0,1 \right)\)，所以 \(P\left( X<1 \right)=P\left( \frac{X-1}{2}<0 \right)=\Phi\left( 0 \right)=0.5\)；同理 \(Y\sim N\left( 2,2^{2} \right)\) 时 \(P\left( Y<2 \right)=0.5\)，用 \(\le\) 同样得 \(0.5\)，所以 \(P\left( X<1 \right)<P\left( Y<2 \right)\) 不成立，故 B错误. 因为样本相关系数的绝对值越接近 \(1\)，成对数据的线性相关程度越强，越接近 \(0\) 则越弱，所以 C，D正确. 故说法错误的是 B.
\end{solution}

\begin{problem}
记 \(S_n\) 为数列 \(\{a_n\}\) 的前 \(n\) 项和. 已知 \(S_n=-n^2+8n\)，则数列 \(\{|a_n|\}\) 的前 12 项和为（\quad）

\choices
    {112}
    {48}
    {80}
    {64}
\end{problem}

\begin{answer}
C
\end{answer}

\begin{solution}
按本页现稿选项与原PDF选项编排不同，但 \(80\) 对应两版 C，结论均为 C. 因为 \(S_{n}=-n^{2}+8n\)，所以 \(a_{1}=S_{1}=7\)；当 \(n\ge 2\) 时 \(a_{n}=S_{n}-S_{n-1}=-2n+9\)，而 \(n=1\) 代入同样得 \(7\)，故对一切 \(n\) 有 \(a_{n}=-2n+9\). 因为 \(a_{n}\ge 0\) 得 \(n\le 4\)，\(a_{n}\le 0\) 得 \(n\ge 5\)（其中 \(a_{4}=1\)，\(a_{5}=-1\)），所以 \(\sum_{k=1}^{12}\left| a_{k} \right|=\sum_{k=1}^{4}a_{k}-\sum_{k=5}^{12}a_{k}=2S_{4}-S_{12}\). 因为 \(S_{4}=16\)，\(S_{12}=-48\)，所以上式为 \(2\times 16-\left( -48 \right)=80\). 故选 C，故排除其他三项.
\end{solution}

\begin{problem}
函数 \(f(x)=0.3^x-\sqrt{x}\) 的零点所在区间是（\quad）

\choices
    {\((0,0.3)\)}
    {\((0.3,0.5)\)}
    {\((0.5,1)\)}
    {\((1,2)\)}
\end{problem}

\begin{answer}
B
\end{answer}

\begin{solution}
因为 \(y=0.3^{x}\) 在 \(\R\) 上单调递减，\(y=\sqrt{x}\) 在 \(\left[ 0,+\infty \right)\) 上单调递增，所以 \(f\left( x \right)=0.3^{x}-\sqrt{x}\) 在定义域 \(\left[ 0,+\infty \right)\) 上单调递减. 因为 \(f\left( 0 \right)=1>0\)，\(0.3<0.5\) 蕴含 \(0.3^{0.3}>0.3^{0.5}=\sqrt{0.3}\)，所以 \(f\left( 0.3 \right)=0.3^{0.3}-\sqrt{0.3}>0\)；因为 \(0.3<0.5\) 蕴含 \(0.3^{0.5}<0.5^{0.5}\)，所以 \(f\left( 0.5 \right)=0.3^{0.5}-\sqrt{0.5}<0\). 因为 \(f\left( x \right)\) 在 \(\left[ 0.3,0.5 \right]\) 上连续且 \(f\left( 0.3 \right)>0\)，\(f\left( 0.5 \right)<0\)，所以由零点存在性定理，存在 \(x_{0}\in\left( 0.3,0.5 \right)\) 使 \(f\left( x_{0} \right)=0\)；又 \(f\left( x \right)\) 单调递减，所以零点唯一，只能落在该区间内. 因为 \(f\left( 0 \right)>0\)，\(f\left( 0.3 \right)>0\)，所以 \(\left( 0,0.3 \right)\) 内无零点，故排除 A；因为 \(f\left( 0.5 \right)<0\) 且递减，所以 \(\left( 0.5,1 \right)\)，\(\left( 1,2 \right)\) 内全为负，故排除 C，D. 故选 B.
\end{solution}

\begin{problem}
\(f(x)=\sin(\omega x+\varphi)\)（\(\omega>0\)，\(-\pi<\varphi<\pi\)），在 \(\left[-\frac{5\pi}{12},\frac{\pi}{12}\right]\) 上单调递增，且 \(x=\frac{\pi}{12}\) 是它的一条对称轴，\((\frac{\pi}{3},0)\) 是它的一个对称中心.

当 \(x\in[0,\frac{\pi}{2}]\) 时，\(f(x)\) 的最小值为（\quad）

\choices
    {\(-\frac{\sqrt3}{2}\)}
    {\(-\frac12\)}
    {1}
    {0}
\end{problem}

\begin{answer}
A
\end{answer}

\begin{solution}
按本页现稿与原PDF题面条件等价，结论均为 A. 设最小正周期为 \(T\)，则 \(T=\frac{2\pi}{\omega}\). 因为 \(f\left( x \right)\) 在 \(\left[ -\frac{5\pi}{12},\frac{\pi}{12} \right]\) 上单调递增，该区间长度为 \(\frac{\pi}{2}\)，而正弦型函数任一单调区间长度不超过 \(\frac{T}{2}\)，所以 \(\frac{\pi}{2}\le \frac{T}{2}\)，即 \(T\ge \pi\)，故 \(0<\omega \le 2\). 因为 \(x=\frac{\pi}{12}\) 是对称轴，所以 \(\omega \cdot \frac{\pi}{12}+\varphi =\frac{\pi}{2}+k\pi\)（\(k\in \Z\)）；因为 \(\left( \frac{\pi}{3},0 \right)\) 是对称中心，所以 \(\omega \cdot \frac{\pi}{3}+\varphi =m\pi\)（\(m\in \Z\)）. 两式相减得 \(\frac{\omega \pi}{4}=\left( m-k \right)\pi -\frac{\pi}{2}\)，即 \(\omega =4\left( m-k \right)-2\)，记 \(n=m-k-1\in \Z\)，则 \(\omega =4n+2\). 反向看，对称轴与对称中心距离为 \(\frac{\pi}{3}-\frac{\pi}{12}=\frac{\pi}{4}\)，必为 \(\frac{\left( 2t+1 \right)T}{4}\)（\(t=0,1,2,\cdots\)），所以 \(\frac{\pi}{4}=\frac{\left( 2t+1 \right)T}{4}\)，即 \(T=\frac{\pi}{2t+1}\). 结合 \(T\ge \pi\) 得 \(2t+1\le 1\)，故 \(t=0\)，于是 \(T=\pi\)，\(\omega =\frac{2\pi}{T}=2\)，这与 \(0<\omega \le 2\) 及 \(\omega =4n+2\) 联立同样只剩 \(n=0\)，\(\omega =2\). 代入得 \(\frac{\pi}{6}+\varphi =\frac{\pi}{2}+k\pi\)，\(\frac{2\pi}{3}+\varphi =m\pi\)，即 \(\varphi =\frac{\pi}{3}+k\pi =m\pi -\frac{2\pi}{3}\). 因为 \(-\pi <\varphi <\pi\)，所以候选 \(\varphi =\frac{\pi}{3}\)（\(k=0\)，\(m=1\)）或 \(\varphi =-\frac{2\pi}{3}\)（\(k=-1\)，\(m=0\)）. 若 \(\varphi =-\frac{2\pi}{3}\)，则 \(2x+\varphi \in \left[ -\frac{3\pi}{2},-\frac{\pi}{2} \right]\) 当 \(x\in \left[ -\frac{5\pi}{12},\frac{\pi}{12} \right]\)，正弦在该相位区间上单调递减，与题设递增矛盾，故舍去；若 \(\varphi =\frac{\pi}{3}\)，则相位为 \(\left[ -\frac{\pi}{2},\frac{\pi}{2} \right]\)，单调递增，符合题设. 故 \(f\left( x \right)=\sin \left( 2x+\frac{\pi}{3} \right)\). 当 \(x\in \left[ 0,\frac{\pi}{2} \right]\) 时 \(2x+\frac{\pi}{3}\in \left[ \frac{\pi}{3},\frac{4\pi}{3} \right]\)，在 \(\left[ \frac{\pi}{3},\frac{\pi}{2} \right]\) 上递增，在 \(\left[ \frac{\pi}{2},\frac{4\pi}{3} \right]\) 上递减，所以最小值在端点取得. 因为 \(\sin \frac{\pi}{3}=\frac{\sqrt{3}}{2}\)，\(\sin \frac{4\pi}{3}=-\frac{\sqrt{3}}{2}\)，所以最小值为 \(-\frac{\sqrt{3}}{2}\). 按本页现稿 C为 \(1\)，按原PDF C为 \(-1\)，两种 C均与最小值不符而排除，故两版结论均为 A，另排除 B，D.
\end{solution}

\begin{problem}
双曲线 \(\frac{x^2}{a^2}-\frac{y^2}{b^2}=1\)（\(a>0\)，\(b>0\)）的左，右焦点分别为 \(F_1\)，\(F_2\)，以右焦点 \(F_2\) 为焦点的抛物线 \(y^2=2px\)（\(p>0\)）与双曲线在第一象限交于 \(P\)，若 \(|PF_1|+|PF_2|=3|F_1F_2|\)，则双曲线离心率 \(e=\)（\quad）

\choices
    {2}
    {5}
    {\(\frac{\sqrt2+1}{2}\)}
    {\(\frac{\sqrt5+1}{2}\)}
\end{problem}

\begin{answer}
A
\end{answer}

\begin{solution}
记双曲线半焦距为 \(c\)，则 \(F_{1}=(-c,0)\)，\(F_{2}=(c,0)\)，\(|F_{1}F_{2}|=2c\)，\(c^{2}=a^{2}+b^{2}\). 因为 \(P\) 在第一象限，在双曲线右支上，所以 \(|PF_{1}|-|PF_{2}|=2a\). 又 \(|PF_{1}|+|PF_{2}|=3|F_{1}F_{2}|=6c\)，两式相加减解得 \(|PF_{1}|=3c+a\)，\(|PF_{2}|=3c-a\). 因为抛物线 \(y^{2}=2px\) 的焦点为 \(\left(\frac{p}{2},0\right)\)，该焦点就是 \(F_{2}(c,0)\)，所以 \(\frac{p}{2}=c\)，即 \(p=2c\)，抛物线为 \(y^{2}=4cx\)，准线为 \(x=-c\). 设 \(P=(x_{0},y_{0})\)（\(x_{0}>0\)，\(y_{0}>0\)），由抛物线定义，\(|PF_{2}|\) 等于 \(P\) 到准线 \(x=-c\) 的距离，所以 \(|PF_{2}|=x_{0}+c\)，故 \(x_{0}=|PF_{2}|-c=2c-a\). 将 \(P\) 代入两条曲线方程，得 \(y_{0}^{2}=4cx_{0}=4c(2c-a)\)，\(\frac{x_{0}^{2}}{a^{2}}-\frac{y_{0}^{2}}{b^{2}}=1\). 代入 \(x_{0}=2c-a\) 得 \(\frac{(2c-a)^{2}}{a^{2}}-\frac{4c(2c-a)}{b^{2}}=1\)，移项得 \(\frac{4c(c-a)}{a^{2}}=\frac{4c(2c-a)}{b^{2}}\)，即 \(\frac{c-a}{a^{2}}=\frac{2c-a}{b^{2}}\). 因为 \(b^{2}=c^{2}-a^{2}\)，所以 \((c^{2}-a^{2})(c-a)=a^{2}(2c-a)\). 记 \(e=\frac{c}{a}>1\)，上式除以 \(a^{3}\) 得 \((e^{2}-1)(e-1)=2e-1\)，即 \(e^{3}-e^{2}-3e+2=0\)，分解得 \((e-2)(e^{2}+e-1)=0\). 方程 \(e^{2}+e-1=0\) 的正根为 \(\frac{\sqrt{5}-1}{2}<1\)，舍去，所以 \(e=2\). 回代检验 \(x_{0}=2c-a=a(2e-1)=3a>0\)，\(y_{0}^{2}=4c(2c-a)>0\)，\(|PF_{2}|=3c-a=5a>0\)，均成立，故选：A.
\end{solution}
\section{填空题}

\begin{problem}
已知 \(\i\) 是虚数单位，则 \(\left|\frac{3+\i}{\i}\right|=\fillinblank{}\).
\end{problem}

\begin{answer}
\(\sqrt{10}\)
\end{answer}

\begin{solution}
因为 \(\frac{3-\i}{\i}=\frac{3}{\i}-1\)，\(\frac{1}{\i}=-\i\)，所以 \(\frac{3-\i}{\i}=-1-3\i\). 故 \(\left|\frac{3-\i}{\i}\right|=\sqrt{(-1)^{2}+(-3)^{2}}=\sqrt{10}\). 若按现稿 \(3+\i\) 计算，\(\frac{3+\i}{\i}=1-3\i\)，模同样为 \(\sqrt{10}\).
\end{solution}

\begin{problem}
在 \((x-1)^6\) 的展开式中，\(x^3\) 项系数为 \(\fillinblank{}\).
\end{problem}

\begin{answer}
\(-20\)
\end{answer}

\begin{solution}
由二项式定理，\((x-1)^{6}=\sum_{k=0}^{6}\mathrm C_{6}^{k}x^{k}(-1)^{6-k}\). 含 \(x^{3}\) 的项对应 \(k=3\)，系数为 \(\mathrm C_{6}^{3}(-1)^{3}=-20\).
\end{solution}

\begin{problem}
\(l_1:x-y+6=0\) 与 \(x\) 轴，\(y\) 轴交点分别为 \(A\)，\(B\)，又与 \((x+1)^2+(y-3)^2=r^2\) 交于 \(C\)，\(D\)，且 \(|AB|=3|CD|\)，则 \(r=\fillinblank{}\).
\end{problem}

\begin{answer}
\(2\)
\end{answer}

\begin{solution}
在 \(l_{1}:x-y+6=0\) 中令 \(y=0\) 得 \(x=-6\)，令 \(x=0\) 得 \(y=6\)，所以 \(A=(-6,0)\)，\(B=(0,6)\)，\(|AB|=\sqrt{6^{2}+6^{2}}=6\sqrt{2}\). 圆心为 \((-1,3)\)，它到 \(l_{1}\) 的距离为 \(d=\frac{|-1-3+6|}{\sqrt{2}}=\sqrt{2}\). 按原PDF \(r>0\)，现稿未显写该条件但同样按半径 \(r>0\) 求解. 设弦 \(CD\) 的中点为 \(M\)，则 \(|CD|=2\sqrt{r^{2}-d^{2}}=2\sqrt{r^{2}-2}\)，其中要求 \(r>d\). 由 \(|AB|=3|CD|\) 得 \(6\sqrt{2}=6\sqrt{r^{2}-2}\)，解得 \(r^{2}=4\). 因为 \(r>0\)，所以 \(r=2\)，满足 \(r>d\)，直线确实交圆于两点.
\end{solution}

\begin{problem}
小桐一周跑 \(2\) 次，一次 \(5\) 圈或 \(6\) 圈. 第一次跑 \(5\) 圈或 \(6\) 圈概率各为 \(0.5\)；若第一次跑 \(5\) 圈，则第二次分别跑 \(5\) 圈，\(6\) 圈的概率为 \(0.4\)，\(0.6\)；若第一次跑 \(6\) 圈，则第二次分别跑 \(5\) 圈，\(6\) 圈的概率为 \(0.6\)，\(0.4\).
则一周跑 \(11\) 圈的概率为 \(\fillinblank{}\)；若一周至少跑 \(11\) 圈为达标，则连续跑 \(4\) 周，记合格周数为 \(X\)，则 \(E(X)=\fillinblank{}\).
\end{problem}

\begin{answer}
\(0.6\)；\(3.2\)
\end{answer}

\begin{solution}
记 \(B\) 为第一次跑 \(5\) 圈的事件，\(C\) 为第二次跑 \(5\) 圈的事件，则 \(P(B)=0.5\)，\(P(C\mid B)=0.4\)，\(P(C\mid B^{\mathrm c})=0.6\). 一周跑 \(11\) 圈当且仅当两次恰为 \(5\) 圈与 \(6\) 圈的组合，由全概率公式，\(P(11)=P(B)P(C^{\mathrm c}\mid B)+P(B^{\mathrm c})P(C\mid B^{\mathrm c})=0.5\times 0.6+0.5\times 0.6=0.6\). 一周跑 \(12\) 圈当且仅当两次都跑 \(6\) 圈，所以 \(P(12)=P(B^{\mathrm c})P(C^{\mathrm c}\mid B^{\mathrm c})=0.5\times 0.4=0.2\). 因此一周至少跑 \(11\) 圈的概率为 \(P(D)=P(11)+P(12)=0.6+0.2=0.8\). 各周跑步圈数相互独立，所以每周围是否达标是独立重复试验，\(X\sim B(4,0.8)\)，故 \(E(X)=4\times 0.8=3.2\).
\end{solution}

\begin{problem}
\(\triangle ABC\) 中，\(D\) 为 \(AB\) 中点，\(\overrightarrow{CE}=\frac13\overrightarrow{CD}\)，\(\overrightarrow{AC}= \bs b\)，\(\overrightarrow{AB}=\bs a\).
则 \(\overrightarrow{AE}=\fillinblank{}\)（用 \(\bs a\)，\(\bs b\) 表示）；若 \(|\overrightarrow{AE}|=5\)，\(\overrightarrow{AE}\perp\overrightarrow{CB}\)，则 \(\overrightarrow{AE}\cdot\overrightarrow{CD}=\fillinblank{}\).

\FigureLayoutDeclare{img/2025/tianjin/q14_fig1.png}{5.5cm}{0bp 0bp 0bp 0bp}
\bitmapfigure[max width=0.42\linewidth]{img/2025/tianjin/q14_fig1.png}
\end{problem}

\begin{answer}
\(\frac16\bs a+\frac23\bs b\)，\(-15\)
\end{answer}

\begin{solution}
以 \(A\) 为起点，记 \(\overrightarrow{AB}=\bs a\)，\(\overrightarrow{AC}=\bs b\). 按本页现稿“\(D\) 为 \(AB\) 中点”与原PDF“\(\overrightarrow{AD}=\frac12\overrightarrow{AB}\)”等价，现稿“\(\overrightarrow{AE}\perp\overrightarrow{CB}\)”与原PDF“\(\overrightarrow{AE}\perp\overrightarrow{BC}\)”为同一直线垂直，两版结论一致. 因为 \(D\) 为 \(AB\) 中点，所以 \(\overrightarrow{AD}=\frac12\bs a\). 因为 \(\overrightarrow{CE}=\frac13\overrightarrow{CD}\)，所以 \(\overrightarrow{AE}=\overrightarrow{AC}+\overrightarrow{CE}=\bs b+\frac13\overrightarrow{CD}\). 又 \(\overrightarrow{CD}=\overrightarrow{AD}-\overrightarrow{AC}=\frac12\bs a-\bs b\)，代入得 \(\overrightarrow{AE}=\bs b+\frac13\left(\frac12\bs a-\bs b\right)=\frac16\bs a+\frac23\bs b\)，这是第一空. 记 \(\bs v=\overrightarrow{AE}\)，则 \(\overrightarrow{BC}=\bs b-\bs a\)，\(\overrightarrow{CD}=\frac12\bs a-\bs b\). 因为 \(\overrightarrow{AE}\perp \overrightarrow{BC}\)，所以 \(\bs v\cdot(\bs b-\bs a)=0\)，即 \(\bs v\cdot\bs a=\bs v\cdot\bs b\)，记此公共值为 \(s\). 于是 \(\bs v\cdot\overrightarrow{CD}=\bs v\cdot\left(\frac12\bs a-\bs b\right)=\frac{s}{2}-s=-\frac{s}{2}\). 另一方面，\(|\bs v|^{2}=\bs v\cdot\bs v=\bs v\cdot\left(\frac16\bs a+\frac23\bs b\right)=\frac{s}{6}+\frac{2s}{3}=\frac{5s}{6}\). 已知 \(|\bs v|=5\)，所以 \(\frac{5s}{6}=25\)，解得 \(s=30\). 故 \(\overrightarrow{AE}\cdot\overrightarrow{CD}=-\frac{s}{2}=-15\)，这是第二空.
\end{solution}

\begin{problem}
若 \(a\)，\(b\in\R\)，对任意 \(x\in[-2,2]\)，都有 \((2a+b)x^2+bx-a-1\le0\)，则 \(2a+b\) 的最小值为 \(\fillinblank{}\).
\end{problem}

\begin{answer}
\(-4\)
\end{answer}

\begin{solution}
记 \(t=2a+b\)，则 \(b=t-2a\)，原不等式化为 \(tx^{2}+(t-2a)x-a-1\le 0\) 对任意 \(x\in[-2,2]\) 成立，即 \(t(x^{2}+x)+a(-2x-1)-1\le 0\). 取 \(x=-\frac12\in[-2,2]\)，此时 \(x^{2}+x=-\frac14\)，\(-2x-1=0\)，不等式变为 \(-\frac{t}{4}-1\le 0\)，解得 \(t\ge -4\). 因此任何满足条件的 \((a,b)\) 必有 \(2a+b\ge -4\). 取 \(a=0\)，\(b=-4\)，则 \(t=-4\)，原不等式左边为 \(-4x^{2}-4x-1=-(2x+1)^{2}\le 0\)，对一切实数 \(x\) 成立，当然对 \(x\in[-2,2]\) 成立，等号在 \(x=-\frac12\) 处取得. 所以 \(-4\) 可以取到，故 \(2a+b\) 的最小值为 \(-4\).
\end{solution}
\section{解答题}

\begin{problem}
在 \(\triangle ABC\) 中，角 \(A\)，\(B\)，\(C\) 对边分别为 \(a\)，\(b\)，\(c\). 已知 \(a\sin B=\sqrt3\,b\cos A\)，\(c-2b=1\)，\(a=\sqrt7\).
\begin{enumerate}
    \item 求 \(A\)；
    \item 求 \(c\)；
    \item 求 \(\sin(A+2B)\).
\end{enumerate}
\end{problem}

\begin{solution}
\begin{enumerate}
\item 因为 \(a\sin B=\sqrt3\,b\cos A\)，由正弦定理 \(\frac{a}{\sin A}=\frac{b}{\sin B}\) 得 \(a\sin B=b\sin A\)，所以 \(b\sin A=\sqrt3\,b\cos A\).因为 \(b>0\)，若 \(\cos A=0\) 则由上式得 \(\sin A=0\)，与 \(\sin^2 A+\cos^2 A=1\) 矛盾，所以 \(\cos A\ne 0\)，故 \(\tan A=\sqrt3\).因为 \(A\in\left(0,\pi\right)\)，所以 \(A=\frac{\pi}{3}\).
\item 因为 \(A=\frac{\pi}{3}\)，所以 \(\cos A=\frac{1}{2}\).由余弦定理 \(a^2=b^2+c^2-2bc\cos A\) 及 \(c-2b=1\) 得 \(c=2b+1\)，代入 \(a=\sqrt7\) 得 \(7=b^2+\left(2b+1\right)^2-b\left(2b+1\right)\)，化简得 \(b^2+b-2=0\)，即 \(\left(b+2\right)\left(b-1\right)=0\).因为 \(b>0\)，所以 \(b=1\)，从而 \(c=3\).经检验 \(1\)，\(3\)，\(\sqrt7\) 满足三角形三边关系，故 \(c=3\).
\item 由余弦定理 \(b^2=a^2+c^2-2ac\cos B\)，代入 \(a=\sqrt7\)，\(b=1\)，\(c=3\) 得 \(1=7+9-6\sqrt7\cos B\)，所以 \(\cos B=\frac{5\sqrt7}{14}\).因为 \(0<\frac{5\sqrt7}{14}<1\)，所以 \(B\in\left(0,\frac{\pi}{2}\right)\) 为锐角，故 \(\sin B>0\)，于是 \(\sin B=\sqrt{1-\cos^2 B}=\sqrt{1-\frac{175}{196}}=\frac{\sqrt{21}}{14}\).从而 \(\cos 2B=2\cos^2 B-1=2\times\frac{175}{196}-1=\frac{11}{14}\)，\(\sin 2B=2\sin B\cos B=2\times\frac{\sqrt{21}}{14}\times\frac{5\sqrt7}{14}=\frac{5\sqrt3}{14}\).因此 \(\sin\left(A+2B\right)=\sin A\cos 2B+\cos A\sin 2B=\frac{\sqrt3}{2}\times\frac{11}{14}+\frac12\times\frac{5\sqrt3}{14}=\frac{4\sqrt3}{7}\).
\end{enumerate}
\end{solution}

\begin{problem}
正方体 \(ABCD-A_1B_1C_1D_1\) 的棱长为 4，\(E\)，\(F\) 分别为 \(A_1D_1\)，\(C_1B_1\) 中点，且 \(CG=3GC_1\).

\begin{enumerate}
\item 求证：\(GF\perp\) 平面 \(FBE\)；
\item 求平面 \(FBE\) 与平面 \(EBG\) 的夹角余弦；
\item 求四面体 \(D-FBE\) 的体积.
\end{enumerate}

\bitmapfigure[max width=0.55\linewidth]{img/2025/tianjin/q17_fig1.png}
\end{problem}

\begin{solution}
\begin{enumerate}
\item 取 \(D\left(0,0,0\right)\)，\(A\left(4,0,0\right)\)，\(B\left(4,4,0\right)\)，\(C\left(0,4,0\right)\)，\(D_1\left(0,0,4\right)\)，\(A_1\left(4,0,4\right)\)，\(B_1\left(4,4,4\right)\)，\(C_1\left(0,4,4\right)\).因为 \(E\)，\(F\) 分别为 \(A_1D_1\)，\(C_1B_1\) 中点，所以 \(E\left(2,0,4\right)\)，\(F\left(2,4,4\right)\).因为 \(G\) 按本页题图在棱 \(CC_1\) 上，与原PDF“\(G\) 在线段 \(CC_1\) 上”一致（现稿文字只写 \(CG=3GC_1\)，未显写该限定），故取线段内分点，而 \(\left|CC_1\right|=4\)，所以 \(\left|CG\right|=3\)，故 \(G\left(0,4,3\right)\).于是 \(\overrightarrow{EF}=\left(0,4,0\right)\)，\(\overrightarrow{EB}=\left(2,4,-4\right)\)，\(\overrightarrow{FG}=\left(-2,0,-1\right)\).设平面 \(FBE\) 的一个法向量为 \(\bs{n}=\left(x,y,z\right)\)，则 \(\overrightarrow{EF}\cdot\bs{n}=4y=0\)，\(\overrightarrow{EB}\cdot\bs{n}=2x+4y-4z=0\).取 \(x=2\) 得 \(y=0\)，\(z=1\)，故可取 \(\bs{n}=\left(2,0,1\right)\).因为 \(\overrightarrow{FG}=\left(-2,0,-1\right)=-\bs{n}\)，所以 \(\overrightarrow{FG}\) 平行于平面 \(FBE\) 的法向量，故 \(GF\perp\) 平面 \(FBE\).
\item 由上问知平面 \(FBE\) 的一个法向量可取 \(\bs{n}_1=\left(2,0,1\right)\).又 \(\overrightarrow{EG}=\left(-2,4,-1\right)\)，\(\overrightarrow{BG}=\left(-4,0,3\right)\).设平面 \(EBG\) 的一个法向量为 \(\bs{n}_2=\left(x,y,z\right)\)，则 \(\overrightarrow{EG}\cdot\bs{n}_2=-2x+4y-z=0\)，\(\overrightarrow{BG}\cdot\bs{n}_2=-4x+3z=0\).取 \(x=6\) 得 \(z=8\)，\(y=5\)，故可取 \(\bs{n}_2=\left(6,5,8\right)\).于是 \(\bs{n}_1\cdot\bs{n}_2=12+0+8=20\)，\(\left|\bs{n}_1\right|=\sqrt{4+0+1}=\sqrt5\)，\(\left|\bs{n}_2\right|=\sqrt{36+25+64}=\sqrt{125}=5\sqrt5\).所以所求夹角余弦为 \(\frac{\left|\bs{n}_1\cdot\bs{n}_2\right|}{\left|\bs{n}_1\right|\left|\bs{n}_2\right|}=\frac{20}{\sqrt5\times 5\sqrt5}=\frac{20}{25}=\frac45\).
\item 因为 \(\overrightarrow{EF}=\left(0,4,0\right)\)，\(\overrightarrow{BF}=\left(-2,0,4\right)\)，所以 \(\overrightarrow{EF}\cdot\overrightarrow{BF}=0\times\left(-2\right)+4\times 0+0\times 4=0\)，故 \(EF\perp BF\).又 \(\left|EF\right|=4\)，\(\left|BF\right|=\sqrt{\left(-2\right)^2+0^2+4^2}=2\sqrt5\)，所以 \(S_{\triangle BFE}=\frac12\left|EF\right|\cdot\left|BF\right|=\frac12\times 4\times 2\sqrt5=4\sqrt5\).又 \(\overrightarrow{DE}=\left(2,0,4\right)\)，因为由第（1）问 \(\overrightarrow{FG}\) 即平面 \(FBE\) 的法向，所以点 \(D\) 到平面 \(FBE\) 的距离可用该法向按投影计算，得 \(d=\frac{\left|\overrightarrow{DE}\cdot\overrightarrow{FG}\right|}{\left|\overrightarrow{FG}\right|}=\frac{\left|2\times\left(-2\right)+0\times 0+4\times\left(-1\right)\right|}{\sqrt5}=\frac{8}{\sqrt5}\).所以 \(V_{D-BFE}=\frac13\cdot d\cdot S_{\triangle BFE}=\frac13\times\frac{8}{\sqrt5}\times 4\sqrt5=\frac{32}{3}\).
\end{enumerate}
\end{solution}

\begin{problem}
已知椭圆 \(\frac{x^2}{a^2}+\frac{y^2}{b^2}=1\)（\(a>b>0\)）的左焦点为 \(F\)，右顶点为 \(A\)，\(P\) 为 \(x=a\) 上一点，且直线 \(PF\) 的斜率为 \(\frac13\)，\(\triangle PFA\) 的面积为 \(\frac32\)，离心率为 \(\frac12\).
\begin{enumerate}
\item 求椭圆方程；
\item 过 \(P\) 的直线与椭圆有唯一交点 \(B\)（异于 \(A\)），证明 \(PF\) 平分 \(\angle AFB\).
\end{enumerate}

\bitmapfigure[max width=0.42\linewidth]{img/2025/tianjin/q18_fig1.png}
\end{problem}

\begin{solution}
\begin{enumerate}
\item 设椭圆半焦距为 \(c\)，则左焦点为 \(F\left(-c,0\right)\)，右顶点为 \(A\left(a,0\right)\).由离心率 \(e=\frac12\) 得 \(\frac{c}{a}=\frac12\)，所以 \(a=2c\).因为 \(P\) 在直线 \(x=a\) 上，设 \(P\left(a,m\right)\)，则直线 \(PF\) 的斜率为 \(\frac{m-0}{a-\left(-c\right)}=\frac{m}{a+c}=\frac13\).因为 \(a+c=3c>0\)，所以 \(m\) 与斜率同号，得 \(m>0\)，即 \(P\) 在 \(x\) 轴上方，且 \(m=c\)，故 \(P\left(2c,c\right)\).又 \(\left|AF\right|=a+c=3c\)，点 \(P\) 到 \(x\) 轴的距离为 \(\left|m\right|=c\)，所以 \(S_{\triangle PFA}=\frac12\times 3c\times c=\frac{3c^2}{2}\).由题意该面积为 \(\frac32\)，得 \(\frac{3c^2}{2}=\frac32\).因为 \(c>0\)，所以 \(c=1\)，从而 \(a=2\)，\(b^2=a^2-c^2=4-1=3\)，故椭圆方程为 \(\frac{x^2}{4}+\frac{y^2}{3}=1\).
\item 由上问知 \(P\left(2,1\right)\)，\(F\left(-1,0\right)\)，\(A\left(2,0\right)\).若过 \(P\) 的直线无斜率，则为 \(x=2\)，与椭圆的唯一公共点为 \(A\left(2,0\right)\)，不符合 \(B\) 异于 \(A\)，所以该直线斜率存在.设其方程为 \(\ell:y=kx+m\).因为 \(P\left(2,1\right)\) 在直线上，所以 \(1=2k+m\)，即 \(m=1-2k\).联立 \(y=kx+m\) 与 \(\frac{x^2}{4}+\frac{y^2}{3}=1\)，消去 \(y\) 得 \(\left(3+4k^2\right)x^2+8kmx+4m^2-12=0\).因为 \(\ell\) 与椭圆只有一个公共点，所以判别式为零，即 \(\left(8km\right)^2-4\left(3+4k^2\right)\left(4m^2-12\right)=0\).两边除以 \(16\) 得 \(4k^2m^2-\left(3+4k^2\right)\left(m^2-3\right)=0\)，展开得 \(-3m^2+9+12k^2=0\)，即 \(4k^2-m^2+3=0\).代入 \(m=1-2k\) 得 \(4k^2-\left(1-2k\right)^2+3=0\)，即 \(4k^2-\left(1-4k+4k^2\right)+3=4k+2=0\)，解得 \(k=-\frac12\)，\(m=2\).故 \(\ell:y=-\frac12x+2\).将其代入椭圆方程得 \(4x^2-8x+4=0\)，即 \(x=1\)，\(y=\frac32\)，故 \(B\left(1,\frac32\right)\)，确异于 \(A\).此时 \(\overrightarrow{FB}=\left(2,\frac32\right)\)，\(\overrightarrow{FP}=\left(3,1\right)\)，\(\overrightarrow{FA}=\left(3,0\right)\).于是 \(\overrightarrow{FB}\cdot\overrightarrow{FP}=6+\frac32=\frac{15}{2}>0\)，\(\overrightarrow{FA}\cdot\overrightarrow{FP}=9>0\)，故 \(\angle BFP\)，\(\angle PFA\) 均为锐角，且 \(\cos\angle BFP=\frac{\overrightarrow{FB}\cdot\overrightarrow{FP}}{\left|\overrightarrow{FB}\right|\left|\overrightarrow{FP}\right|}=\frac{3\sqrt{10}}{10}\)，\(\cos\angle PFA=\frac{\overrightarrow{FA}\cdot\overrightarrow{FP}}{\left|\overrightarrow{FA}\right|\left|\overrightarrow{FP}\right|}=\frac{3\sqrt{10}}{10}\).因为两角都在 \(\left(0,\frac{\pi}{2}\right)\) 内且余弦相等，所以 \(\angle BFP=\angle PFA\)，故 \(PF\) 平分 \(\angle AFB\).
\end{enumerate}
\end{solution}

\begin{problem}
已知数列 \(\{a_n\}\) 为等差数列，\(\{b_n\}\) 为等比数列，\(a_1=b_1=2\)，\(a_2=b_2+1\)，\(a_3=b_3\).
\begin{enumerate}
\item 求 \(\{a_n\}\)，\(\{b_n\}\) 的通项；
\item \(\forall n\in\N^*\)，\(I=\{0,1\}\)，定义 \(T_n=\{p_1a_1b_1+p_2a_2b_2+\cdots+p_na_nb_n\mid p_i\in I,i=1,\dots,n\}\).
\begin{enumerate}
\item 求证：对任意 \(t\in T_n\)，有 \(t<a_{n+1}b_{n+1}\)；
\item 求 \(T_n\) 的所有元素之和.
\end{enumerate}
\end{enumerate}
\end{problem}

\begin{solution}
\begin{enumerate}
\item 设等差数列 \(\left\{a_n\right\}\) 的公差为 \(d\)，等比数列 \(\left\{b_n\right\}\) 的公比为 \(q\).由 \(a_1=b_1=2\)，\(a_2=b_2+1\)，\(a_3=b_3\) 得 \(2+d=2q+1\)，\(2+2d=2q^2\).由第一式得 \(d=2q-1\)，代入第二式得 \(2+2\left(2q-1\right)=2q^2\)，即 \(q^2-2q=0\)，故 \(q=0\) 或 \(q=2\)，对应 \(d=-1\) 或 \(d=3\).若 \(q=0\)，则 \(b_1=2\)，\(b_2=0\)，此时 \(b_3/b_2\) 无意义，\(\left\{b_n\right\}\) 不成等比数列，故舍去.所以 \(d=3\)，\(q=2\).故 \(a_n=2+3\left(n-1\right)=3n-1\)，\(b_n=2\times 2^{n-1}=2^n\).
\item
\begin{enumerate}
\item 由上问得 \(a_nb_n=\left(3n-1\right)2^n>0\).对任意 \(t\in T_n\)，因为每项系数 \(p_i\in\left\{0,1\right\}\)，所以 \(t\le a_1b_1+a_2b_2+\cdots+a_nb_n\).记 \(S_n=a_1b_1+a_2b_2+\cdots+a_nb_n\)，则 \(S_n=2\times 2+5\times 2^2+\cdots+\left(3n-1\right)2^n\).于是 \(2S_n=2\times 2^2+5\times 2^3+\cdots+\left(3n-1\right)2^{n+1}\).两式相减得 \(-S_n=4+3\left(2^2+2^3+\cdots+2^n\right)-\left(3n-1\right)2^{n+1}\).因为 \(2^2+2^3+\cdots+2^n=2^{n+1}-4\)，所以 \(-S_n=4+3\left(2^{n+1}-4\right)-\left(3n-1\right)2^{n+1}=-8+\left(4-3n\right)2^{n+1}\)，故 \(S_n=8+\left(3n-4\right)2^{n+1}\).又 \(a_{n+1}b_{n+1}=\left(3n+2\right)2^{n+1}\)，所以 \(a_{n+1}b_{n+1}-S_n=6\times 2^{n+1}-8\).当 \(n=1\) 时该差为 \(16>0\)，且随 \(n\) 增大而增大，故对一切 \(n\ge 1\) 恒大于零.因此对任意 \(t\in T_n\) 都有 \(t\le S_n<a_{n+1}b_{n+1}\).
\item 在 \(T_n\) 的全部 \(2^n\) 个元素中，固定某一下标 \(i\) 并令 \(p_i=1\)，其余 \(n-1\) 个系数仍可在 \(\left\{0,1\right\}\) 中自由取值，共有 \(2^{n-1}\) 种组合，所以每个 \(a_ib_i\) 在求和时恰出现 \(2^{n-1}\) 次.因此 \(T_n\) 的所有元素之和为 \(2^{n-1}\sum_{i=1}^n a_ib_i=2^{n-1}S_n=2^{n-1}\left[8+\left(3n-4\right)2^{n+1}\right]\).
\end{enumerate}
\end{enumerate}
\end{solution}

\begin{problem}
已知 \(f(x)=ax-(\ln x)^2\).
\begin{enumerate}
\item 当 \(a=1\) 时，求 \(f(x)\) 在点 \((1,f(1))\) 处的切线方程；
\item 当 \(f(x)\) 有三个零点 \(x_1<x_2<x_3\) 时，
\begin{enumerate}
    \item 求 \(a\) 的取值范围；
    \item 证明 \((\ln x_2-\ln x_1)\cdot \ln x_3<\frac{4\e}{\e-1}\).
\end{enumerate}
\end{enumerate}
\end{problem}

\begin{solution}
\begin{enumerate}
\item 当 \(a=1\) 时，\(f\left(x\right)=x-\left(\ln x\right)^2\) 的定义域为 \(x>0\)，求导得 \(f'\left(x\right)=1-\frac{2\ln x}{x}\).所以 \(f'\left(1\right)=1\)，且 \(f\left(1\right)=1\)，故切线为 \(y-1=1\times\left(x-1\right)\)，即 \(y=x\).
\item
\begin{enumerate}
\item 令 \(f\left(x\right)=0\) 得 \(a=\frac{\left(\ln x\right)^2}{x}\).设 \(g\left(x\right)=\frac{\left(\ln x\right)^2}{x}\)，\(x>0\)，则 \(g'\left(x\right)=\frac{2\ln x-\left(\ln x\right)^2}{x^2}=\frac{\ln x\left(2-\ln x\right)}{x^2}\).由 \(g'\left(x\right)=0\) 得 \(x=1\) 或 \(x=\e^2\)，且 \(g\left(1\right)=0\)，\(g\left(\e^2\right)=\frac{4}{\e^2}\).当 \(0<x<1\) 时 \(g'\left(x\right)<0\)，当 \(1<x<\e^2\) 时 \(g'\left(x\right)>0\)，当 \(x>\e^2\) 时 \(g'\left(x\right)<0\).又当 \(x\to 0^+\) 时 \(g\left(x\right)\to +\infty\)，当 \(x\to +\infty\) 时 \(g\left(x\right)\to 0\).
\solutionfigure{\bitmapfigure[max width=0.55\linewidth]{img/2025/tianjin/q20_solution_fig1.png}}
因此 \(g\) 在 \(\left(0,1\right)\) 上从 \(+\infty\) 严格递减到 \(0\)，在 \(\left(1,\e^2\right)\) 上从 \(0\) 严格递增到 \(\frac{4}{\e^2}\)，在 \(\left(\e^2,+\infty\right)\) 上从 \(\frac{4}{\e^2}\) 严格递减到 \(0\).故方程 \(a=g\left(x\right)\) 在 \(\left(0,+\infty\right)\) 内恰有三个根，当且仅当 \(0<a<\frac{4}{\e^2}\).
\item 由上问 \(g\) 的单调性及 \(a=g\left(x\right)\) 有三个根，推出 \(0<x_1<1<x_2<\e^2<x_3\).设 \(\ln x_1=t_1\)，\(\ln x_2=t_2\)，\(\ln x_3=t_3\)，则 \(t_1<0<t_2<2<t_3\)，且 \(a\e^{t_1}=t_1^2\)，\(a\e^{t_2}=t_2^2\)，\(a\e^{t_3}=t_3^2\).由后两式取对数得 \(\ln a+t_2=2\ln t_2\)，\(\ln a+t_3=2\ln t_3\)，两式相减得 \(t_3-t_2=2\left(\ln t_3-\ln t_2\right)\). 记 \(u=\frac{t_3}{t_2}>1\)，\(s=\sqrt{u}>1\)，则 \(\frac{t_3-t_2}{\ln t_3-\ln t_2}=t_2\frac{u-1}{\ln u}\). 因为当 \(s>1\) 时 \(\ln s<\frac12\left(s-\frac1s\right)\)（令 \(h\left(s\right)=\frac12\left(s-\frac1s\right)-\ln s\)，则 \(h\left(1\right)=0\)，\(h'\left(s\right)=\frac{\left(s-1\right)^2}{2s^2}>0\)），所以 \(\ln u=2\ln s<s-\frac1s=\frac{u-1}{\sqrt{u}}\)，故 \(\frac{u-1}{\ln u}>\sqrt{u}\). 于是 \(\frac{t_3-t_2}{\ln t_3-\ln t_2}>t_2\sqrt{u}=\sqrt{t_2t_3}\). 结合左端等于 \(2\) 得 \(2>\sqrt{t_2t_3}\)，所以 \(t_2t_3<4\).所要证的 \(\left(\ln x_2-\ln x_1\right)\ln x_3=\left(t_2-t_1\right)t_3=t_2t_3+\left(-t_1t_3\right)\)，只需分别界定两项之和.因为 \(t_1<0\) 且 \(t_1^2=a\e^{t_1}<a\)，其中 \(\e^{t_1}<1\) 用了 \(t_1<0\)，所以 \(-t_1<\sqrt{a}\).又由 \(a\e^{t_3}=t_3^2\) 得 \(\sqrt{a}\,t_3=\frac{t_3^2}{\e^{\frac{t_3}{2}}}\).于是 \(-t_1t_3<\sqrt{a}\,t_3=\frac{t_3^2}{\e^{\frac{t_3}{2}}}\).设 \(\varphi\left(t\right)=\frac{t^2}{\e^{\frac{t}{2}}}\)，\(t>2\)，则 \(\varphi'\left(t\right)=\frac{t\left(4-t\right)}{2\e^{\frac{t}{2}}}\).所以当 \(2<t<4\) 时 \(\varphi'\left(t\right)>0\)，当 \(t>4\) 时 \(\varphi'\left(t\right)<0\)，故 \(\varphi\left(t\right)\) 在 \(t=4\) 处取得最大值 \(\varphi\left(4\right)=\frac{16}{\e^2}\).而 \(\frac{16}{\e^2}<\frac{4}{\e-1}\)，因为此式等价于 \(4\left(\e-1\right)<\e^2\)，即 \(\e^2-4\e+4=\left(\e-2\right)^2>0\).从而 \(-t_1t_3<\frac{16}{\e^2}<\frac{4}{\e-1}\).结合 \(t_2t_3<4\) 得 \(\left(t_2-t_1\right)t_3<4+\frac{4}{\e-1}=\frac{4\e}{\e-1}\)，即 \(\left(\ln x_2-\ln x_1\right)\cdot\ln x_3<\frac{4\e}{\e-1}\).
\end{enumerate}
\end{enumerate}
\end{solution}
